Consider the following activity diagram:



\begin{centering}\includegraphics[min width=0.4\textwidth=0.6\textwidth,min height=0.4\textwidth=0.6\textwidth,keepaspectratio,max width=0.9\textwidth,max height=0.9\textwidth]{73/Diagram.pdf}\end{centering}

Translate the given activity diagram into a Petri net (on paper or in your head) and then state the total count of nodes (places and transitions),
the count of auxiliary places and the count of auxiliary transitions in the net.

To do this, enter your answer as in the following example:\begin{verbatim}FindAuxiliaryPetriNodesSolution {countOfPetriNodes = 10, countOfAuxiliaryPlaces = 2, countOfAuxiliaryTransitions = 3}\end{verbatim}In this example, the resulting net contains 10 nodes in total, of which 2 are auxiliary places and 3 are auxiliary transitions.

Hint on the translation to a Petri net: For final nodes no additional places are introduced. They are realised in a way that a token is consumed, i.e. disappears from the net at that position. If an additional transition is required to realise this behavior at a position in the diagram where there is a final node, this transition does not count as auxiliary node.

